% !TEX root = ../main.tex
\providecommand{\main}{..}
\documentclass[../main.tex]{subfiles}

\begin{document}
\chapter{Literature Review}
\section{AdaRound}
\subsection{Notation}
Let us establish the notation for variables in Table \ref{tab:adaround_notation}.

\begin{table}[H]
\centering
\caption{Notation for this Literature Review}
\label{tab:adaround_notation}
\begin{tabular}{ll}
\hline
\textbf{Symbol} & \textbf{Description} \\ \hline
$w_j$             & Weights of a layer $j$ \\
$x_j$             & Input to a layer $j$ \\
$b$             & Target bit-width quantization \\
$\text{scale}$  & Continuous scaling factor \\
$\lfloor \cdot \rceil$ & Round-to-nearest-integer operator \\
$\lfloor \cdot \rfloor$ & Floor operator \\
%$\mathbf{V}$  & Matrix of learned AdaRound parameter \\
$\mathbf{W}, \hat{\mathbf{W}}$    & Matrices of unquantized and quantized weights \\
$\mathcal{L}$ & A loss function \\
\hline
\end{tabular}
\end{table}

\subsection{Introduction to AdaRound}
AdaRound~\cite{nagelAdaptiveRoundingPostTraining2020a} is a novel rounding method, based on the key insight that round-to-nearest is suboptimal for post-training quantization (PTQ).

The authors' proposed method is to instead minimize the difference in task loss of the quantized model and the unquantized model. The task loss of layer $j$ of the unquantized model is $\mathcal{L}(x, y, w_j)$, and the task loss of the same layer of the quantized model $\mathcal{L}(x, y, w_j + \Delta w_j)$. Therefore, the authors' objective is the following.

\begin{align}
\underset{\Delta w_j}{\arg\min} \left(\mathcal{L}(x, y, w_j + \Delta w_j) - \mathcal{L}(x, y, w_j)\right)
\end{align}

The authors use a second order Taylor expansion to approximate the above objective.

\begin{align}
\mathcal{L}(x, y, w_j + \Delta w_j) - \mathcal{L}(x, y, w_j) \approx \nabla\mathcal{L}(x, y, w_j)^T \Delta w_j + \frac{1}{2} \Delta w_j^T \mathbf{H} \Delta w_j
\end{align}

Where $\mathbf{H}$ is the Hessian $\nabla^2\mathcal{L}(x, y, w_j)$. The authors then ignore the first term by asserting that the gradient $\nabla\mathcal{L}(x, y, w_j)$ is $0$ for converged pretrained models, which simplifies into the following.

\begin{align}
\mathcal{L}(x, y, w_j + \Delta w_j) - \mathcal{L}(x, y, w_j) \approx \frac{1}{2} \Delta w_j^T \mathbf{H} \Delta w_j
\end{align}

Finally, the authors formulate the optimization of weigh rounding for layer $j$ as the following Quadratic Unconstrained Binary Optimization (QUBO) problem.

\begin{align}
\underset{\Delta w_j}{\arg\min} \Delta w_j^T \mathbf{H} \Delta w_j
\end{align}

The above problem is NP-hard. To solve this efficiently, AdaRound uses the local mean squared error alongside a regularizer as a proxy. The regularizer encourages the optimization variables to converge towards either $0$ or $1$.

\todo{Derivation}

\begin{align}
\arg\min_{\mathbf{V}} \|\mathbf{W}x - \hat{\mathbf{W}}x\|_F^2 + f_\text{reg}(\mathbf{V})
\end{align}

The regularizer uses an annealed parameter $\beta$ and is defined as follows.

\begin{align}
f_{\text{reg}}(\mathbf{V}) = \sum_{i, j} 1 - \lvert 2h(\mathbf{V}_{i,j}) - 1\rvert^\beta
\end{align}

The above objective is then optimized over

\begin{align}
\hat{\mathbf{W}} &= \text{scale} \cdot \text{clamp} \left( \lfloor \frac{\mathbf{W}}{\text{scale}} \rfloor +h(\mathbf{V}), q_\text{min}, q_\text{max}\right), \\ 
q_\text{min} &= -2^{b-1}, \\
q_\text{max} &= 2^{b-1}-1.
\end{align}

The authors define $h$ as a rectified sigmoid as follows, although it can be any differentiable function with a range of $h(\mathbf{V}_{i,j}) \in [0, 1]$.

\begin{align}
h(\mathbf{V}_{i,j}) &= \text{clamp}\left( \sigma(\mathbf{V}_{i, j}) (\zeta - \gamma) + \gamma, 0, 1\right), \\
\zeta &= 1.1, \\
\gamma &= -0.1.
\end{align}

The use of a rectified sigmoid is justified as it has non-vanishing gradients near $0$ and $1$, and helps the learning process as they previously used a regularizer to encourage $h(\mathbf{V}_{i, j})$ to move to either $0$ or $1$.

\subsection{AdaRound Algorithm}
We implement AdaRound in Algorithm \ref{alg:adaround}.

\begin{algorithm}[H]
\caption{AdaRound}\label{alg:adaround}
\begin{algorithmic}[1]
\Require $\mathbf{W}$, $X$ (calibration data), $t$ (number of iterations), $b$ (target bit width), $\beta$ (regularization parameter)
\Ensure Quantized model $\mathbf{W}_q$, scales

\State $x_{in} \gets \textsc{CacheActivations}(\mathbf{W}, X)$ \Comment{Forward pass on calibration data to cache input activations}
\For{$w_j$ in $\mathbf{W}$}
\State $q_{\text{min}} \gets -2^{b - 1}$
\State $q_{\text{max}} \gets 2^{b - 1} - 1$
\State $\text{scales}_j \gets \frac{\max(w_j)}{q_{\text{max}}}$ \Comment{Store scale for output}
\State $\mathbf{V}_j \gets h^{-1}(w_j - \lfloor w_j \rfloor)$ \Comment{Initialize $\mathbf{V}_j$ so that it corresponds to the fraction of each weight}

\For{$i \in \{0 \ldots t\}$}
\State $\beta \gets \textsc{Anneal}(\beta)$
\State $\hat{w}_j \gets \text{clamp}\left(\lfloor \frac{w_j}{\text{scales}_j} \rfloor + h(\mathbf{V}_j), q_\text{min}, q_\text{max}\right)$
\State $\mathcal{L} \gets \|w_j x_{in} - \hat{w}_j x_{in}\|_F^2 + \sum_{i, j} \left(1 - \lvert 2h(\mathbf{V}_{i,j}) - 1\rvert^\beta\right)$
\State $\mathbf{V} \gets \mathbf{V} - \nabla \mathcal{L}$ \Comment{via autograd}
\EndFor
\State $\mathbf{W}_{Q_{i, j}} \gets \begin{cases}
    \lceil w_{i, j} \rceil, &\text{where } h(\mathbf{V}_{i,j}) \geq 0.5 \\
    \lfloor w_{i, j} \rfloor, &\text{otherwise}
\end{cases}$ \Comment{Finalize quantization}
\EndFor
\State \Return $\mathbf{W}_Q$, scales


\end{algorithmic}
\end{algorithm}

\subsection{Empirical Results}
\todo{Report AdaRound results}

\subsection{Limitations}
\todo{Two things here: 1) zero gradient assumption, 2) data leakage}

\section{Inversion Method}
\todo{Check if OK to directly copy from previous paper}
In this section, we will report our novel inversion method (also previously published as \cite{liuAdaRoundInversion}) to reconstruct representations of the calibration data used by AdaRound, given access to a model before and after AdaRound quantization.

\subsection{Formulation of Leakage Channel}
AdaRound's objective function is expected to approach zero after training. Let us define the weight perturbation matrix as $\Delta W = W - \hat{W}$. The objective function can be rewritten as the following.
\begin{align}
\|Wx - \hat{W}x\|_F^2 = \|\Delta W x\|_F^2
\end{align}
After optimization, we expect the following.
\begin{align}
\|\Delta W x\|_F^2 \approx 0 \iff \Delta W x \approx 0
\end{align}

We notice that $x \subset \mathcal{N}(\Delta W)$. In other words, the training data $x$ used to calibrate the AdaRound process lies within the null space of the weight perturbation matrix $\Delta W$. 

\subsection{Method}
An attacker with access to the original weights $W$ and the quantized weights $\hat{W}$ can generate a sample $\hat{x}$ that lies within this null space. We can initialize a $\hat{x}$ from random noise and apply Stochastic Gradient Descent with the objective $\arg\min_{\hat{x}} \|\Delta W \hat{x}\|_2^2$, as described in Algorithm \ref{alg:inversion}.

\begin{algorithm}[H]
\caption{Proposed Inversion Method}\label{alg:inversion}
\begin{algorithmic}[1]
\Require $W$, $\hat{W}$, $\hat{x}_{\text{init}}$, Number of Iterations $K$, Learning Rate $\eta$
\State $\Delta W \gets W - \hat{W}$
\State $\hat{x} \gets \hat{x}_\text{init}$
\For{$k = 1$ \textbf{to} $K$}
    \State $\mathcal{L}  \gets \sum \left( \Delta W  \tanh(\hat{x}) \right)^2$ (Forward pass)
    \State $\frac{\partial \mathcal{L}}{\partial \hat{x}} \gets \text{Autograd}(\mathcal{L}, \hat{x})$ (Backward pass via autograd~\cite{paszkePyTorchImperativeStyle2019})
    \State $\hat{x} \gets \hat{x} - \eta \frac{\partial\mathcal{L}}{\partial{\hat{x}}}$ (Update $\hat{x}$)
\EndFor
\State \Return $\tanh(\hat{x})$
\end{algorithmic}
\end{algorithm}

Note that the objective can be rewritten as the following.
\begin{align}
\|\Delta W \hat{x}\|_2^2 = (\Delta W \hat{x})^T(\Delta W \hat{x}) = \hat{x}^T (\Delta W^T \Delta W) \hat{x}
\end{align}
A simple optimization can therefore be made by computing the Gram matrix $G = \Delta W^T \Delta W$ once, then use it to calculate the objective as $\hat{x}^TG\hat{x}$ in each iteration. (See Algorithm \ref{alg:inversion_optimized})

\begin{algorithm}[H]
\caption{Proposed Inversion with Optimization}\label{alg:inversion_optimized}
\begin{algorithmic}[1]
\Require $W$, $\hat{W}$, $\hat{x}_{\text{init}}$, Number of Iterations $K$, Learning Rate $\eta$
\State $\Delta W \gets W - \hat{W}$
\State $\hat{x} \gets \hat{x}_\text{init}$
\State $G \gets \Delta W^T \Delta W$
\For{$k = 1$ \textbf{to} $K$}
    \State $\mathcal{L}  \gets \sum \left( \tanh(\hat{x})^TG  \tanh(\hat{x}) \right)$ (Forward pass)
    \State $\frac{\partial \mathcal{L}}{\partial \hat{x}} \gets \text{Autograd}(\mathcal{L}, \hat{x})$ (Backward pass via autograd~\cite{paszkePyTorchImperativeStyle2019})
    \State $\hat{x} \gets \hat{x} - \eta \frac{\partial\mathcal{L}}{\partial{\hat{x}}}$ (Update $\hat{x}$)
\EndFor
\State \Return $\tanh(\hat{x})$
\end{algorithmic}
\end{algorithm}

\subsection{Evaluation}

We trained MLPs on these datasets: Olivetti Faces~\cite{JMLR:v12:pedregosa11a}, Chest X-Ray Images (Pneumonia)~\cite{kermanyIdentifyingMedicalDiagnoses2018}, and the Weizmann Horse Database~\cite{weizmannHorseDB}, on a single computer (CPU: Intel Xeon w9-3495X, GPU: 1x NVIDIA GeForce RTX 4090). We perform 500 iterations of AdaRound quantization using the Adam optimizer, then obtain $\Delta W$, the difference between the quantized and the original model. For each model, we generate 10 random samples from a Gaussian distribution, then train for 3000 iterations using the Adam optimizer. We report the SSIM (Structural Similarity Index Measure), PSNR (Peak Signal-to-Noise ratio), and Normalized RMSE (Root Mean Squared Error) between the mean of the calibration data and the mean of 10 generated samples before and after our inversion method in Table \ref{tab:adaround_inversion_metrics}. We also report a visualization of the mean calibration sample vs. the mean generated sample in Figures \ref{fig:adaround_inversion_olivetti}, \ref{fig:adaround_inversion_xray}, and \ref{fig:adaround_inversion_horse}.

\begin{table}[H]
\centering
\caption{Metrics from Random Initialization vs. Our Proposed Method}
\label{tab:adaround_inversion_metrics}
\begin{tabular}{|l|l|l|l|}
\hline
                        & SSIM $\uparrow$       & PSNR $\uparrow$        & NRMSE $\downarrow$    \\ \hline
Olivetti Faces (Before) & 0.00                  & 3.50                   & 1.16                  \\ \hline
Olivetti Faces (After)  & \textbf{0.14 (+0.14)} & \textbf{12.83 (+9.33)} & \textbf{0.40 (-0.76)} \\ \hline
Chest X-Ray (Before)    & 0.00                  & 2.61                   & 1.10                  \\ \hline
Chest X-Ray (After)     & \textbf{0.07 (+0.07)} & \textbf{10.13 (+7.52)} & \textbf{0.46 (-0.64)} \\ \hline
Horse DB (Before)       & 0.00                  & 2.60                   & 1.11                  \\ \hline
Horse DB (After)        & \textbf{0.07 (+0.07)} & \textbf{10.15 (+7.55)} & \textbf{0.47 (-0.64)} \\ \hline
\end{tabular}
\end{table}

\begin{figure}[H]
    \centering
    \includegraphics[width=0.6\textwidth]{img/adaround_inversion/16acd81804674b3b8d69cbe4807a3e0f-result.png}
    \caption{Results on Olivetti Faces}
    \label{fig:adaround_inversion_olivetti}
\end{figure}

\begin{figure}[H]
    \centering
    \includegraphics[width=0.6\textwidth]{img/adaround_inversion/04a80ceb974248ffb6864725da78e909-result.png}
    \caption{Results on Chest X-Ray Images (Pneumonia)}
    \label{fig:adaround_inversion_xray}
\end{figure}

\begin{figure}[H]
    \centering
    \includegraphics[width=0.6\textwidth]{img/adaround_inversion/8bbfd491fcb644dda0d9a7b3db367b66-result.png}
    \caption{Results on Weizmann Horse Database}
    \label{fig:adaround_inversion_horse}
\end{figure}

We notice an increase in SSIM and PSNR, and a decrease in Normalized RMSE, implying that the mean of the generated samples resembles the mean of the calibration samples to some extent.

\subsection{Ablation Study}
We perform an ablation study by applying our inversion method to round-to-nearest (RTN) quantization. Since RTN does not use calibration data, we compare our method's samples to the mean of the training data instead. We report the same metrics as the main experiment (SSIM, PSNR, and Normalized RMSE) in Table \ref{tab:adaround_inversion_ablation} and a visualization in Figures \ref{fig:olivetti_ablation}, \ref{fig:xray_ablation}, and \ref{fig:horse_ablation}.

\begin{table}[H]
\centering
\caption{Results from Ablation Study}
\label{tab:adaround_inversion_ablation}
\begin{tabular}{|l|l|l|l|}
\hline
 &  SSIM $\uparrow$ & PSNR $\uparrow$ & NRMSE $\downarrow$ \\ \hline
Olivetti Faces (Before) & 0.00 & 2.89 & 1.12 \\ \hline
Olivetti Faces (After) & 0.00 & 3.90 (+1.01) & 1.00 (-0.12) \\ \hline
Chest X-Ray (Before) & 0.00 & 2.77 & 1.11 \\ \hline
Chest X-Ray (After) & 0.00 & 3.64 (+0.87) & 1.00 (-0.11) \\ \hline
Horse DB (Before) & 0.00 & 2.57 & 1.10 \\ \hline
Horse DB (After) & 0.00 & 3.38 (+0.81) & 1.00 (-0.10) \\ \hline
\end{tabular}
\end{table}

\begin{figure}[H]
    \centering
    \includegraphics[width=0.6\textwidth]{img/adaround_inversion/ablation-31bebceaa4bb4cd2b9305c23af7e5d8c-result.png}
    \caption{Ablation Study Results on Olivetti Faces}
    \label{fig:olivetti_ablation}
\end{figure}

\begin{figure}[H]
    \centering
    \includegraphics[width=0.6\textwidth]{img/adaround_inversion/ablation-3f5b592012394a26a9ef69d4e24cf68f-result.png}
    \caption{Ablation Study Results on Chest X-Ray Images (Pneumonia)}
    \label{fig:xray_ablation}
\end{figure}

\begin{figure}[H]
    \centering
    \includegraphics[width=0.6\textwidth]{img/adaround_inversion/ablation-bdc15cdff0e94745bebd2a5af058167e-result.png}
    \caption{Ablation Study Results on Weizmann Horse Database}
    \label{fig:horse_ablation}
\end{figure}

We notice 0.00 SSIM on all datasets, implying that the generated samples do not look like the calibration data at all.

\subsection{Discussion}
\subsubsection{Main Experiment}
From the metrics, we observe an increase in SSIM and PSNR, and a decrease in Normalized RMSE after training. In the images, we can see the silhouettes and outlines of the calibration set, although it is quite difficult to make out. From the results, we can conclude that the $\Delta W$ of a model quantized using AdaRound does contain information about the calibration data. 

\subsubsection{Ablation Study}
Interestingly, we see a slight improvement in PSNR and NRMSE in the ablation study. However, crucially the SSIM shows that round-to-nearest quantization is not vulnerable to our method. In other words, our method indeed extracts information about the calibration data in data-dependent quantization methods.

\subsubsection{Overall}
We demonstrate that AdaRound encodes information from the calibration data into the weight perturbation matrix. By targeting the null space of this matrix, our proposed method generates samples that resembles outlines of the calibration data. While the reconstructed data only represents the mean of the calibration data rather than exact instances, the leakage is significant enough to warrant reconsideration of the security of data-dependent PTQ methods. 

\end{document}